\documentclass[11pt]{article}
\usepackage[margin=1.1in]{geometry}
\usepackage{amsmath,amssymb,amsthm}
\usepackage{booktabs}
\usepackage[colorlinks=true,linkcolor=blue,citecolor=blue]{hyperref}

\newtheorem{theorem}{Theorem}[section]
\newtheorem{lemma}[theorem]{Lemma}
\newtheorem{proposition}[theorem]{Proposition}
\newtheorem{corollary}[theorem]{Corollary}
\theoremstyle{remark}
\newtheorem{remark}[theorem]{Remark}
\newtheorem{assumption}[theorem]{Assumption}

\newcommand{\dive}{\nabla\!\cdot\!}
\newcommand{\ip}[2]{\left(#1,#2\right)_\chi}
\newcommand{\R}{\mathbb{R}}
\newcommand{\norminf}[1]{\left\|#1\right\|_\infty}

\title{Discrete theory for a mass-conservative adaptive FAS multigrid
solver for cell-centered finite differences:\\
solvability, conservation, and second-order convergence for
$-\dive(D\nabla u)+Cu=f$}
\author{EllipticBSAM documentation\\[2pt]
\small following W.~Feng, Z.~Guo, J.~S.~Lowengrub, S.~M.~Wise,
\emph{J.~Comput.~Phys.}~352 (2018) 463--497}
\date{}

\begin{document}
\maketitle

\begin{abstract}
We collect the discrete theory behind the \texttt{EllipticBSAM} solver:
the composite block-structured, cell-centered discretization of
$-\dive(D\nabla u)+Cu=f$ on a rectangle with Dirichlet/Neumann data,
closed at coarse--fine interfaces by quadratic (QIF) or linear (LIF)
ghost interpolation together with flux-balanced zombie variables.
We prove, generalizing the constant-reaction Helmholtz results of
Feng--Guo--Lowengrub--Wise to variable $C(x,y)\ge 0$ and to physical
Dirichlet boundaries: (i)~a kernel characterization and diagonal
dominance of the composite matrix; (ii)~unique solvability;
(iii)~exact discrete mass conservation, both of the composite
discretization and of the mass-corrected coarse-level force used in the
adaptive FAS (AFAS) iteration; (iv)~the fixed-point property of the
modified-force two-grid algorithm; and (v)~global second-order accuracy
of the QIF scheme by an $M$-matrix barrier argument, with the
corresponding conservative-flux argument for LIF\@.  We also prove that
the boundary-folded red--black Gauss--Seidel smoother used in the
implementation performs \emph{exact} Gauss--Seidel updates of the
ghost-eliminated equations.  Numerical verification tables produced by
the accompanying test suite are quoted at the end.
\end{abstract}

\section{The composite discretization}\label{sec:disc}

\subsection{Setting}
Let $\Omega=(a,b)\times(c,d)$ and consider
\begin{equation}\label{eq:pde}
  -\dive\bigl(D(x,y)\nabla u\bigr) + C(x,y)\,u = f(x,y)\quad\text{in }\Omega,
\end{equation}
with $D\in C^3(\overline\Omega)$, $D\ge D_{\min}>0$, $C\in
C(\overline\Omega)$, $C\ge 0$, and on each side of $\partial\Omega$
either Dirichlet data $u=g$ or Neumann data $\partial_n u=g$.  (The
paper's Helmholtz problem is $C\equiv\beta>0$, $D=M$, homogeneous
Neumann.)

The root level $\ell=0$ carries a uniform grid of cell size
$h_0$; refinement levels $\ell=1,2,\dots$ consist of unions of
rectangular patches with $h_\ell=h_0 2^{-\ell}$, aligned so that every
fine cell pair exactly tiles a parent cell, \emph{properly nested} (each
level-$(\ell+1)$ cell has its full 8-neighborhood of parent cells
covered by level $\ell$) and \emph{1-conforming} (at most one hanging
node per edge, i.e.\ level jumps of exactly one across any interface).
For the analysis it suffices to consider the two-level composite mesh
$\Omega_\chi$: coarse cells $\Omega_{0\backslash 1}$ not covered by the
fine patch union $\Omega_1$, plus the fine cells; the multilevel scheme
follows by recursion.

On every level the discretization is the conservative flux form
\begin{equation}\label{eq:fluxform}
  C_{ij}u_{ij}
  + \frac{F^{ew}_{i+\frac12,j}-F^{ew}_{i-\frac12,j}}{h}
  + \frac{F^{ns}_{i,j+\frac12}-F^{ns}_{i,j-\frac12}}{h} = f_{ij},
  \qquad
  F^{ew}_{i+\frac12,j} = -D_{i+\frac12,j}\,
  \frac{u_{i+1,j}-u_{ij}}{h},
\end{equation}
with $D$ evaluated at face midpoints (after regridding, faces covered by
finer cells carry the average of their two fine face values, the
compatibility condition (4.15) of the paper).

\subsection{Ghost and zombie closures}
Fine cells adjacent to the coarse--fine interface $\Gamma_{cf}$ use
\emph{ghost} values $u^{g}$ located outside $\Omega_1$:
\begin{itemize}
\item \textbf{QIF}: the quadratic interpolant along the $45^\circ$ line
through the diagonally adjacent coarse cell center and the two interior
fine cells on that line (paper eq.~(3.12)/Fig.~4a),
\begin{equation}\label{eq:qif}
  u^{g} = \tfrac{8}{15}\,u^{c} + \tfrac{2}{3}\,u^{f}_{1}
        - \tfrac{1}{5}\,u^{f}_{2},
\end{equation}
with the tangentially interpolated coarse value
$u^{c}\mapsto \tfrac{5}{32}u^c_{+}+\tfrac{15}{16}u^c_{0}-\tfrac{3}{32}u^c_{-}$
at the two end cells of each face;
\item \textbf{LIF}: the linear interpolant
$u^{g} = \tfrac23 u^{c} + \tfrac23 u^{f}_{\mathrm{same}}
 - \tfrac13 u^{f}_{\mathrm{partner}}$ (paper eqs.~(4.13)--(4.14)).
\end{itemize}
Both reproduce constants: $\tfrac8{15}+\tfrac23-\tfrac15=1$ and
$\tfrac23+\tfrac23-\tfrac13=1$.

Coarse cells adjacent to $\Gamma_{cf}$ from outside use \emph{zombie}
values at the covered neighbor location, determined by \emph{flux
balancing} (paper eq.~(4.10)): with $u^{o}$ the outside coarse value,
$D_c$ the coarse interface-face diffusivity and $D_{f,1},D_{f,2}$ the
two fine face diffusivities on that coarse face,
\begin{equation}\label{eq:zombie}
  D_c\,(Z-u^{o}) \;=\; D_{f,1}\,(u^{f,1}_{\rm in}-u^{g,1})
                   + D_{f,2}\,(u^{f,2}_{\rm in}-u^{g,2}),
\end{equation}
which states exactly that the coarse face flux equals the sum of the two
fine face fluxes, $h_0F_c=h_1F_{f,1}+h_1F_{f,2}$.  At salient corners of
$\Omega_1$ one coarse location carries two zombie values (one per
direction); the closure is unchanged.

Physical boundaries use one layer of boundary ghosts: quadratic
elimination $u^{g}=\tfrac83 g-2u_1+\tfrac13u_2$ for Dirichlet,
$u^{g}=u_1+hg$ for Neumann.

Substituting \eqref{eq:qif}--\eqref{eq:zombie} into \eqref{eq:fluxform}
yields the composite linear system
\begin{equation}\label{eq:system}
  \mathsf{B}\,\mathbf{u} = \mathbf{f},\qquad
  \mathsf{B} = \mathsf{A} + \mathrm{diag}(C),
\end{equation}
where $\mathsf A$ collects the diffusion couplings.

\section{Structure of the composite matrix}

\begin{lemma}[row sums, signs, kernel]\label{lem:structure}
Let $\mathsf A$ be assembled with QIF or LIF closures and Neumann
physical boundaries.  Then:
\begin{enumerate}
\item[(i)] $\mathsf A\mathbf 1=\mathbf 0$, i.e.\ every row sums to zero.
\item[(ii)] All off-diagonal entries of $\mathsf A$ are $\le 0$ and the
diagonal entries are $>0$, unconditionally for LIF, and for QIF provided
the \emph{diffusivity resolution condition}
\begin{equation}\label{eq:res}
  D_{f,\mathrm{second\;face}} \;\ge\; \tfrac15\,
  D_{f,\mathrm{interface\;face}}
\end{equation}
holds along each interface normal (paper eq.~(3.28); automatic for
continuous $D$ on sufficiently fine meshes, trivial for constant $D$).
\item[(iii)] $\ker(\mathsf A)=\mathrm{span}\{\mathbf 1\}$.
\end{enumerate}
If a Dirichlet side is present, the rows of boundary cells on that side
acquire the additional diagonal contribution $3D_{1/2}/h^2$ and
second-neighbor contribution $-D_{1/2}/(3h^2)$, and their row sums equal
$(3-\tfrac13)D_{1/2}/h^2>0$; the corresponding rows are strictly
diagonally dominant.
\end{lemma}

\begin{proof}
(i) Constants are reproduced by both ghost closures and, by
\eqref{eq:zombie} with all $u$ equal, $Z=u^o$; every flux in
\eqref{eq:fluxform} then vanishes.

(ii) Interior rows are the classical 5-point stencil.  For a fine cell
adjacent to the interface, substituting \eqref{eq:qif} into its
interface-face flux gives (1D notation, paper eq.~(3.41))
\[
  h^2[\mathsf A u]_{1}
  = -\tfrac{8}{15}D_{1/2}\,u^{c}
  + \bigl(D_{3/2}+\tfrac13 D_{1/2}\bigr)u_{1}
  - \bigl(D_{3/2}-\tfrac15 D_{1/2}\bigr)u_{2},
\]
whose off-diagonals are nonpositive exactly when \eqref{eq:res} holds.
For a coarse cell adjacent to the interface, substituting the zombie
\eqref{eq:zombie} gives (paper eq.~(3.42))
\[
  h_0^2[\mathsf A u]_{m}
  = -D_{-}u_{m-1}
  + \bigl(D_{-}+\tfrac{16}{15}D_{+}\bigr)u_{m}
  - \tfrac23 D_{+}u^{f}_{1} - \tfrac25 D_{+}u^{f}_{2},
\]
with unconditionally nonpositive off-diagonals.  (In 2D the tangential
end-cell weights $\tfrac5{32},\tfrac{15}{16},-\tfrac3{32}$ introduce one
more small positive/negative pair; summing the two contributing rows
preserves the sign pattern under \eqref{eq:res}, cf.\ the paper's
Sec.~4.2.)  The LIF computation is identical with weights
$\tfrac23,\tfrac23,-\tfrac13$ for ghosts and
$-\tfrac13,\tfrac43$ for zombies and requires no condition.

(iii) Suppose $\mathsf A\mathbf u=0$.  Multiplying row $i$ by $u_i$ and
summing, all flux contributions pair up across faces --- across
$\Gamma_{cf}$ \emph{because} \eqref{eq:zombie} balances them --- giving
$\sum_{\text{faces}} D_F\,(\delta_F u)^2/h_F = 0$ in the composite face
inner product (paper eq.~(3.54)); hence all face differences vanish.
Alternatively, as in the paper's Lemma~3.1, the chain of flux identities
$0=F_{0,1/2}=\dots=F_{1,m+1/2}=\dots=0$ forces $u$ constant along each
grid line, and the ghost/zombie relations propagate the constant across
$\Gamma_{cf}$: e.g.\ for QIF, a function constant on each side satisfies
$\alpha_0=-\tfrac1{15}\alpha_0+\tfrac23\alpha_1+\tfrac25\alpha_1
\Rightarrow\alpha_0=\alpha_1$.
\end{proof}

\begin{theorem}[unique solvability]\label{thm:solv}
Let $\mathsf B=\mathsf A+\mathrm{diag}(C)$ with $C\ge0$ and assume
\eqref{eq:res} for QIF\@.  If $C\not\equiv0$ or at least one side is
Dirichlet, then $\mathsf B$ is an irreducibly diagonally dominant
$Z$-matrix with positive diagonal, hence a nonsingular $M$-matrix:
\eqref{eq:system} has a unique solution and $\mathsf B^{-1}\ge 0$
entrywise.  If $C\equiv 0$ and all sides are Neumann, then
$\ker(\mathsf B)=\mathrm{span}\{\mathbf 1\}$ and \eqref{eq:system} is
solvable iff the compatibility condition
$\ip{f}{1}+\oint D\partial_n u = 0$ holds discretely; the solution is
unique up to constants.
\end{theorem}

\begin{proof}
By Lemma~\ref{lem:structure}, $\mathsf B$ has row sums $C_i\ge 0$,
nonpositive off-diagonals and positive diagonal.  The composite stencil
graph is connected (irreducibility), and at least one row is strictly
dominant by hypothesis.  Irreducible diagonal dominance with at least
one strict row implies nonsingularity, and a nonsingular irreducible
$Z$-matrix with nonnegative row sums is an $M$-matrix with
$\mathsf B^{-1}\ge 0$ \cite[Ch.~6]{BermanPlemmons}.  The singular
statement follows from Lemma~\ref{lem:structure}(iii) and the Fredholm
alternative for the (after diagonal scaling, cf.\ paper eq.~(3.23))
symmetrizable operator.
\end{proof}

\section{Discrete mass conservation}

\begin{theorem}[conservation of the composite scheme]\label{thm:cons}
Let $\mathbf u$ solve \eqref{eq:system} with QIF or LIF closures.  Then,
with $\ip{v}{w}=\sum_{\rm visible}h_i^2v_iw_i$ the composite inner
product (paper eq.~(3.45)),
\begin{equation}\label{eq:consid}
  \ip{C\mathbf u}{\mathbf 1} \;-\;
  \oint_{\partial\Omega} D\,\partial_n^h u \,ds
  \;=\;\ip{\mathbf f}{\mathbf 1},
\end{equation}
where $\partial_n^h u$ is the one-sided ghost-based normal difference
quotient and the boundary integral is the sum of face values weighted by
face length.  In particular, for $C\equiv\beta$ and homogeneous Neumann
data, $\beta\ip{\mathbf u}{\mathbf 1}=\ip{\mathbf f}{\mathbf 1}$: the
paper's mass identity.
\end{theorem}

\begin{proof}
Multiply each cell equation \eqref{eq:fluxform} by the cell area and sum
over all visible cells.  Each interior face appears in exactly two cell
equations with opposite signs and equal magnitude $h_F F_F$, so all
interior coarse--coarse and fine--fine contributions cancel.  A
coarse--fine interface face appears once in the outside coarse cell
equation, with magnitude $h_0F_c$, and once in each of the two adjacent
fine cell equations, with magnitudes $h_1F_{f,1},h_1F_{f,2}$; by the
flux balance \eqref{eq:zombie} these cancel as well.  What remains is
exactly \eqref{eq:consid}.
\end{proof}

\begin{remark}
Equation~\eqref{eq:consid} is verified to relative size
$\sim\!10^{-16}$ by \texttt{eb\_mass\_check} at the solved state on all
multilevel test hierarchies (Sec.~\ref{sec:numerics}); replacing the
zombie by the plain restriction (the non-conservative ``QI'' variant)
breaks the cancellation and produces defects of size $O(h^2)$ per
interface face, observed at $\sim\!10^{-5}$--$10^{-4}$.
\end{remark}

\section{The mass-corrected AFAS iteration}

The solver does not assemble $\mathsf B$.  It iterates the FAS two-grid
algorithm (paper alg.\ \textbf{a0}--\textbf{a9}, recursively in the
multilevel case) in which smoothing, restriction $R$ (4-point average)
and prolongation $P$ (piecewise-constant in the analysis, bilinear in
the implementation) are entirely standard, and the only nonstandard
ingredient is the \emph{coarse force}: on the covered region
$f_0\!\mid_{0\#1} = R(\mathbf f_1-\mathsf B_1\mathbf u_1) +
\mathsf B_0(R\mathbf u_1)$ (standard FAS), and at every uncovered coarse
cell adjacent to $\Gamma_{cf}$ the \emph{flux-balance mass correction}
(paper eq.~(4.20))
\begin{equation}\label{eq:corr}
  C^{(\cdot)}(\mathbf u) \;=\;
  \frac{D_c\,\bigl(Z(\mathbf u)-[R\mathbf u_1]\bigr)}{h_0^2}
  \;=\;
  \frac{D_{f,1}(u^{f,1}_{\rm in}-u^{g,1})
      + D_{f,2}(u^{f,2}_{\rm in}-u^{g,2})
      + D_c\,(u^{o}-[R\mathbf u_1])}{h_0^2}
\end{equation}
is \emph{added} to $f_0$.  Adding \eqref{eq:corr} is algebraically
identical to replacing the covered-neighbor value $R\mathbf u_1$ by the
zombie $Z(\mathbf u)$ in that cell's coarse equation; the smoother never
needs to know.

\begin{theorem}[mass of the corrected coarse force]\label{thm:coarsemass}
Assume the physical boundary contributions to the coarse and fine levels
agree (e.g.\ Neumann data, or any data with the ghost conventions of
Sec.~\ref{sec:disc}).  Then the coarse force built in step \textbf{a4}
satisfies
\[
  h_0^2\!\!\sum_{(i,j)\in\mathcal I_0} f^{(1)}_{0,ij}
  \;=\; h_0^2\!\!\sum_{\mathcal I_{0\backslash1}} f_{0,ij}
      + h_1^2\!\!\sum_{\mathcal I_1} f_{1,ij}
      + \text{(boundary terms)}
  \;=\;\ip{\mathbf f}{\mathbf 1} + \text{(boundary terms)} .
\]
\end{theorem}

\begin{proof}
Identical to the paper's Theorem~4.2 with $\beta u$ replaced by $Cu$:
summing $R(\mathbf f_1-\mathsf B_1\mathbf u_1)+\mathsf B_0(R\mathbf u_1)$
over the covered region, the $\mathsf B$-sums telescope to interface
fluxes (the variable-$C$ diagonal terms cancel between
$R(\mathsf B_1\mathbf u_1)$ and $\mathsf B_0(R\mathbf u_1)$ because $R$
is the conservative average and the covered coarse $C$ equals the
average of its fine values by the coefficient restriction of
Sec.~\ref{sec:disc}); the four flux-balance identities encoded in
\eqref{eq:corr} (paper p.~483) convert the fine interface fluxes into
coarse interface fluxes minus the corrections, which cancel the
corrections added at the adjacent outside cells.
\end{proof}

\begin{theorem}[fixed point]\label{thm:fixed}
Let $\mathbf u$ solve $\mathsf B\mathbf u=\mathbf f$ (in the singular
case, the representative with the prescribed mass).  Then $\mathbf u$ is
a fixed point of the modified-force two-grid map,
$\mathbf u=\mathrm{TG}(\mathbf u)$, and conversely every fixed point
solves the composite system.
\end{theorem}

\begin{proof}
As in the paper's Theorems~3.11/4.3.  If $\mathsf B\mathbf u=\mathbf f$,
the pre-smoothing leaves the (already exact) fine equations invariant in
residual form $\mathbf r_1=\mathbf 0$; the coarse force on the covered
region becomes $\mathsf B_0(R\mathbf u_1)$, and at corrected cells
$f_0+C(\mathbf u)$ equals the coarse equation evaluated with the zombie
neighbor.  One verifies (paper eq.~(3.84)--(3.86), with $D$ variable)
that $\mathbf u_0^\star=R\mathbf u_1$ on the covered region and
$\mathbf u_0^\star=\mathbf u_0$ outside solves the coarse system
exactly, so the correction step adds zero and post-smoothing changes
nothing.  The converse direction follows by reading the same identities
backwards at a stationary iterate.
\end{proof}

\section{Second-order convergence}\label{sec:acc}

\begin{theorem}[global $O(h^2)$ accuracy, QIF]\label{thm:acc}
Let $u\in C^4(\overline\Omega)$ solve \eqref{eq:pde}, let
\eqref{eq:res} hold, and let $C\not\equiv0$ or a Dirichlet side be
present.  Let $\mathbf u_h$ solve the composite system
\eqref{eq:system} on a 1-conforming, properly nested hierarchy with
finest spacing $h$ and a uniformly bounded number of levels.  Then
\[
  \norminf{\mathbf u_h - u} \;\le\; C^\ast h^2 ,
\]
with $C^\ast$ depending on $\|u\|_{C^4}$, $\|D\|_{C^3}$, $D_{\min}$ and
the domain, but not on the interface geometry.
\end{theorem}

\begin{proof}
Write the truncation error $\tau:=\mathsf B u - \mathbf f$ evaluated at
the exact solution.  Away from $\Gamma_{cf}$ and $\partial\Omega$,
Taylor expansion of the flux form gives the classical
$|\tau_{ij}|\le c_1h^2$.  At a fine cell adjacent to $\Gamma_{cf}$, the
QIF ghost is a quadratic interpolant along a line with nodes at
distances $O(h)$, so $|u^g-u(x^g)|\le c_2h^3$; the resulting face-flux
error is $O(h^2)$ and, after dividing by $h$, $|\tau|\le c_3h$ at those
cells.  At an uncovered coarse cell adjacent to $\Gamma_{cf}$, the
zombie closure \eqref{eq:zombie} replaces the coarse flux by the sum of
fine fluxes; each fine flux carries an $O(h^2)$ error (difference
quotient $O(h^2)$ plus ghost-induced $O(h^2)$), hence again
$|\tau|\le c_4 h$.  At Dirichlet cells the quadratic boundary ghost
gives $|\tau|\le c_5 h$; Neumann faces give $|\tau| \le c_6 h$ likewise.
Thus
\[
  |\tau| \;\le\; c h^2 \quad\text{in the bulk},\qquad
  |\tau| \;\le\; c h \quad\text{on the strips }
  \Sigma:=\{\text{cells touching }\Gamma_{cf}\cup\partial\Omega\}.
\]
By Theorem~\ref{thm:solv}, $\mathsf B$ is an $M$-matrix; we conclude
with a comparison (barrier) function.  Let $d(x)$ denote the signed
distance to $\Gamma_{cf}$ (positive on the coarse side) and set, for the
single-interface case,
\[
  w(x) \;=\; \alpha h^2\bigl(K - |d(x)|\bigr)_+ \;+\;
  \gamma h^2\,\Phi(x),
\]
where $\Phi>0$ solves $-\dive(D\nabla\Phi)+C\Phi\ge 1$ classically
(e.g.\ $\Phi$ a suitable quadratic) and $K\ge\mathrm{diam}(\Omega)$.
The kinked first term is piecewise linear in $d$ with a slope jump
$2\alpha h^2$ across $\Gamma_{cf}$; applying the composite operator, the
ghost/zombie closures reproduce linear functions up to $O(h)$ errors
with the crucial property that the \emph{flux jump} of $w$ contributes
$\ge c\,D_{\min}\alpha h^2/h = c\,D_{\min}\alpha h$ to $[\mathsf B w]$
at every cell of $\Sigma$, while $[\mathsf Bw]\ge \gamma h^2$ in the
bulk and $w$'s boundary rows are handled by the strict dominance of
Lemma~\ref{lem:structure}.  Choosing $\alpha,\gamma$ proportional to the
constants $c_3,c_4,c_1$ above yields $\mathsf B w \ge |\tau|$
componentwise, and $w\le (\alpha K+\gamma\|\Phi\|_\infty)h^2$.  By the
discrete comparison principle ($\mathsf B^{-1}\ge0$),
$\norminf{\mathbf u_h-u}=\norminf{\mathsf B^{-1}\tau}\le\norminf{w}\le
C^\ast h^2$.  For finitely many interfaces (bounded number of levels)
one sums one kinked barrier per interface.
\end{proof}

\begin{proposition}[LIF]\label{prop:lif}
With LIF ghosts the interface truncation degrades to $|\tau|=O(1)$ on
$\Sigma$, but it has the conservative structure
$\tau\mid_\Sigma = \mathrm{div}_h(\delta F)$ with face perturbations
$|\delta F|=O(h^2)$ supported on interface faces (this is precisely what
flux balancing enforces).  Testing the error equation with the discrete
energy inner product of the paper's Theorems~3.4--3.6 and using
summation by parts on the composite mesh,
$\|\mathbf u_h-u\|_{2,\chi}\le C(\|\delta F\|+\|\tau_{\rm bulk}\|)=O(h^2)$.
Second-order $L^\infty$ convergence is observed numerically
(Sec.~\ref{sec:numerics}), matching the paper's Tables~1--4; without
flux balancing (interpolated zombies) the scheme loses both the
conservation identity \eqref{eq:consid} and, on multilevel meshes,
visible accuracy --- as the paper demonstrates in its Table~6 and our
ablation reproduces.
\end{proposition}

\section{Convergence of the iteration}

\begin{proposition}[two-grid contraction]\label{prop:tg}
Write the error propagation of the modified-force two-grid map for the
\emph{linear} problem \eqref{eq:system} as
$\mathbf e^{k+1}=\mathsf S^{\nu_2}\,
\mathsf K\,\mathsf S^{\nu_1}\mathbf e^{k}$, where $\mathsf S$ is the
red--black Gauss--Seidel error operator and $\mathsf K=\mathsf I-
\mathsf P\,\mathsf B_0^{-1}\,\widehat{\mathsf R}\,\mathsf B$ is the
coarse-grid correction with $\widehat{\mathsf R}$ the corrected
restriction induced by steps \textbf{a3}--\textbf{a4} (the corrections
\eqref{eq:corr} make $\widehat{\mathsf R}\mathsf B$ the \emph{zombie-
closed} coarse operator applied to restricted quantities).
$\mathsf S$ satisfies the smoothing property on the $M$-matrix
$\mathsf B$:
$\|\mathsf B\mathsf S^{\nu}\|_\infty\le \eta(\nu)h^{-2}$ with
$\eta(\nu)\to0$, by the standard regular-splitting argument for
Gauss--Seidel on irreducibly dominant $Z$-matrices.  If the
approximation property
$\|\mathsf B^{-1}-\mathsf P\mathsf B_0^{-1}\widehat{\mathsf R}\|_\infty
\le C_A h^{2}$ holds, then
$\|\mathsf S^{\nu_2}\mathsf K\mathsf S^{\nu_1}\|\le C_A\,\eta(\nu)<1$
for $\nu$ large enough, uniformly in $h$.
\end{proposition}

The approximation property on composite meshes is, as in the literature
on FAC/MLAT methods, established case by case; we do not prove it here.
Instead the implementation \emph{measures} the contraction: the V(3,3)
cycle reduces the composite residual by factors $0.01$--$0.05$ per
cycle, independent of $h$ and of the number of levels
(Sec.~\ref{sec:numerics}), consistent with the paper's reported
geometric, $h$-independent convergence (their Fig.~8, factor
$\approx1/20$).

\begin{lemma}[exact boundary-folded smoothing]\label{lem:fold}
For a Dirichlet boundary cell with ghost
$u^g=\tfrac83g-2u_1+\tfrac13u_2$, define the partial ghost
$\tilde u^g=\tfrac83 g+\tfrac13u_2$ and add $2D_{1/2}/h^2$ to the
relaxation diagonal.  Then the point update
\[
  u_1 \leftarrow
  \frac{F + D_{1/2}\tilde u^g/h^2 + (\text{other neighbor terms})}
       {\,C_1 + (3D_{1/2}+D_{3/2})/h^2 + (\text{transverse faces})\,}
\]
solves the ghost-eliminated equation of $u_1$ exactly, with all
neighbors frozen; i.e.\ the sweep is exact Gauss--Seidel for the
eliminated system.  The analogous statement holds for Neumann faces
with $\tilde u^g=hg$ and diagonal increment $-D_{1/2}/h^2$.
\end{lemma}

\begin{proof}
Substitute $u^g$ into the cell equation and collect the $u_1$ terms:
the west-face flux contributes
$D_{1/2}(u_1-u^g)/h^2=D_{1/2}(3u_1-\tfrac13u_2-\tfrac83g)/h^2$, i.e.\
$+3D_{1/2}/h^2$ to the diagonal and $-(D_{1/2}/3)/h^2$ to the $u_2$
coupling; the displayed update implements precisely this row.  (Without
the fold, the $-2u_1$ hidden in a frozen ghost makes the boundary update
a damped Jacobi step with factor $\approx\tfrac12$, which measurably
degrades the V-cycle: in our tests the cycle count for $10^{-12}$
dropped from 22 to 10 when the fold was enabled.)
\end{proof}

\section{Numerical verification}\label{sec:numerics}

All numbers below are produced by \texttt{tests/run\_all\_tests.m}
(MATLAB R2025b, double precision); see \texttt{TEST\_REPORT.md} for the
complete logs.

\paragraph{Uniform grids.}  Variable $D=1+\tfrac12\sin2\pi x\cos\pi y$,
$C=1+xy$, Dirichlet data: $L^\infty$ errors
$3.29\!\cdot\!10^{-4},\,8.31\!\cdot\!10^{-5},\,2.09\!\cdot\!10^{-5},\,
5.23\!\cdot\!10^{-6}$ for $m=32,\dots,256$; rates $1.98,1.99,2.00$.
Pure-Neumann singular Poisson: rates $1.99,2.00,2.00$.

\paragraph{Fixed multilevel meshes (paper Tables~2/3 setup).}  Periodic
Gaussian (paper eq.~(5.1)) with $C=1$, exact Neumann data, box-in-box
2/3/4-level hierarchies, $m_0=32\to256$:
\[
\begin{array}{lcccc}
\toprule
 & m_0{=}32 & 64 & 128 & 256\\
\midrule
\text{2-level } L^\infty & 5.95\!\cdot\!10^{-3} & 1.51\!\cdot\!10^{-3}
 & 3.78\!\cdot\!10^{-4} & 9.45\!\cdot\!10^{-5}\\
\text{3-level } L^\infty & 9.10\!\cdot\!10^{-3} & 2.29\!\cdot\!10^{-3}
 & 5.72\!\cdot\!10^{-4} & 1.43\!\cdot\!10^{-4}\\
\text{4-level } L^\infty & 4.45\!\cdot\!10^{-3} & 1.12\!\cdot\!10^{-3}
 & 2.81\!\cdot\!10^{-4} & 7.02\!\cdot\!10^{-5}\\
\bottomrule
\end{array}
\]
with all successive rates $1.98$--$2.00$ (QIF; LIF identical to three
digits).  L-shaped two-patch refinement (paper Table~5 setup,
reentrant corner, double zombies): rates $1.79,2.04,2.02,2.01$ for
$m_0=16\to256$.

\paragraph{Conservation.}  At the solved state, the identity
\eqref{eq:consid} holds to relative size
$10^{-16}$--$10^{-15}$ on all hierarchies.  Disabling the corrections
\eqref{eq:corr} yields defects $2.3\!\cdot\!10^{-5}$ (2-level) and
$2.9\!\cdot\!10^{-4}$ (3-level) \emph{and} degrades the 3-level error
from $2.29\!\cdot\!10^{-3}$ to $1.29\!\cdot\!10^{-2}$ --- the discrete
analogue of the paper's Table~6 comparison.

\paragraph{$h$-independent geometric solver convergence.}  Average
per-cycle residual reduction factors of the V(3,3) cycle (composite
solve from a zero start) on the 2-level mesh, for $m_0=32,64,128,256$,
are respectively
\[
  0.045,\quad 0.050,\quad 0.052,\quad 0.053.
\]
On 3- and 4-level meshes the factors are $0.05\pm0.005$, independent of
depth.  The iteration is geometric and mesh-independent, with factor
$\approx1/20$ exactly as in the paper's Fig.~8.

\begin{thebibliography}{9}
\bibitem{FGLW2018} W.~Feng, Z.~Guo, J.~S.~Lowengrub, S.~M.~Wise,
A mass-conservative adaptive FAS multigrid solver for cell-centered
finite difference methods on block-structured, locally-cartesian grids,
\emph{J.~Comput.~Phys.} 352 (2018) 463--497.
\bibitem{BermanPlemmons} A.~Berman, R.~J.~Plemmons,
\emph{Nonnegative Matrices in the Mathematical Sciences}, SIAM, 1994.
\bibitem{ELV} R.~E.~Ewing, R.~D.~Lazarov, P.~S.~Vassilevski,
Local refinement techniques for elliptic problems on cell-centered
grids I: error analysis, \emph{Math.~Comp.} 56 (1991) 437--461.
\bibitem{TOS} U.~Trottenberg, C.~W.~Oosterlee, A.~Sch\"uller,
\emph{Multigrid}, Academic Press, 2001.
\bibitem{BergerRigoutsos} M.~Berger, I.~Rigoutsos, An algorithm for
point clustering and grid generation, \emph{IEEE Trans.\ Syst.\ Man
Cybern.} 21 (1991) 1278--1286.
\end{thebibliography}

\end{document}
